% Generated draft: requires mathematical and source review.
\documentclass[11pt]{article}
\usepackage[margin=25mm]{geometry}
\usepackage{fontspec}
\setmainfont{Latin Modern Roman}
\usepackage{amsmath,amssymb,mathtools,unicode-math}
\setmathfont{Latin Modern Math}
\usepackage{enumitem,needspace,xurl,hyperref}
\setlength{\parindent}{0pt}
\setlength{\parskip}{5pt}
\setlength{\emergencystretch}{3em}
\newcommand{\N}{\mathbb{N}}
\newcommand{\Z}{\mathbb{Z}}
\newcommand{\Q}{\mathbb{Q}}
\newcommand{\R}{\mathbb{R}}
\newcommand{\C}{\mathbb{C}}
\newcommand{\F}{\mathbb{F}}
\newcommand{\A}{\mathbb{A}}
\newcommand{\PP}{\mathbb{P}}
\newcommand{\E}{\mathbb{E}}
\newcommand{\eps}{\varepsilon}
\newcommand{\dd}{\,\mathrm{d}}
\begin{document}
\section*{UnsolvedMath ID: 20002303}
\textbf{Source problem code:} catalogue entry
\subsection{Definitions and mathematical statement}
Work in a fixed computable countable signature. A computable structure has a computable atomic diagram on a computable domain. Use the standard tuple/back-and-forth convention for Scott rank \(\operatorname{SR}(A)\), the supremum of the tuple orbit-isolation ranks plus one, and let \(\omega_1^{\mathrm{CK}}\) be the least noncomputable ordinal.

The supplied aid reconstructs the workshop's computable-approximability question as follows: for every computable infinitary sentence \(\varphi\) true in \(A\), must there be a computable \(B\models\varphi\) with \(\operatorname{SR}(B)<\omega_1^{\mathrm{CK}}\)? Computable infinitary sentences allow effectively enumerated countable conjunctions/disjunctions and finite-string quantification in \(L_{\omega_1,\omega}\). This is low-Scott-rank approximation of each individual sentence, not elementary equivalence to a low-rank model or the stronger uniform coding notion sometimes called strong computable approximability. The principal extract omits the preceding full definition, so its rank convention must be checked against the cited workshop/theorem.

The aid reports a published negative answer due to Matthew Harrison-Trainor. For each
\[
\alpha\in\{\omega_1^{\mathrm{CK}},\ \omega_1^{\mathrm{CK}}+1\},
\]
there is a computable structure \(A_\alpha\) and a computable \(\Pi^c_2\) sentence \(\psi_\alpha\) such that
\[
A_\alpha\models\psi_\alpha,\qquad
\forall B\text{ countable}\quad B\models\psi_\alpha
\Rightarrow\operatorname{SR}(B)=\alpha.
\]
Here \(\Pi^c_2\) is the computable infinitary complexity class, not a claim that the sentence is a finitary first-order formula. Such a sentence forbids every low-rank computable approximation, and therefore refutes the reconstructed universal question. The exact theorem citation and rank normalization are not supplied.

The aid adds a finite-marker amplification. Fix one barrier pair \((A,\psi)\) at either displayed high rank. In a fixed tagged two-sorted language retain the old structure only on the main sort, and put only equality on a disjoint marker sort. Let \(\psi^P\) be the relativization of \(\psi\) to the main sort, without fixing the marker size. For each positive integer \(k\), let
\[
C_k=(A,F_k),\qquad |F_k|=k.
\]
The reported conclusion is that the \(C_k\) are computable, pairwise nonisomorphic, all satisfy the same computable barrier sentence \(\psi^P\), all have exactly \(\operatorname{SR}(C_k)=\operatorname{SR}(A)\), and all fail computable approximability in the stated sense. Tags and absence of cross-sort relations/functions are part of the construction; the finite pure marker cannot create new high-rank orbit complexity. This amplification builds on the cited high-rank existence theorem and does not prove that theorem independently.
\subsection{Short English statement}
Must every computable structure's true computable-infinitary sentences have low-Scott-rank computable models? The supplied aid reports a published negative answer using a single high-rank barrier sentence. Adding a finite tagged marker produces infinitely many nonisomorphic counterexamples sharing one barrier.
\subsection{Sources}
\begin{itemize}
\item Calvert, subsection 1.1, \emph{Computable approximability}, Question 1, AIM \emph{Computable Stability Theory Problem Session}, dated 12 August 2013: \url{https://aimath.org/pastworkshops/computestabproblems.pdf}.
\item Matthew Harrison-Trainor's published high-Scott-rank barrier theorem, as supplied by the research aid; no title, year, identifier or complete convention statement is provided.
\item Local registry \texttt{data/unsolvedmath/problems.json}, permanent ID 20002303, code \texttt{catalogue entry}; canonical \texttt{aim-logic-notes.json}, \texttt{notes[78]}; updated 14 August 2026. Its verification reports prior PDF page-1 inspection separating the date/next-page header and line-break artifacts from the actual question.
\item Attempt 1: counterexample, problem status at run refuted, rigor label proved, novelty confidence low; current solved and solved-in-literature labels. The machine-generated aid is not peer reviewed and requires independent verification before citation.
\end{itemize}
\end{document}
